% Soluzioni della simulazione d'esame (Appendice A.2 della dispensa)
\chapter{Soluzioni -- Simulazione d'esame}

\section*{Prova 1 -- Quiz a crocette}

Conta le risposte giuste: con 16 o più sei in buona forma; sotto le 12, ripassa i moduli
delle domande sbagliate (indicati tra parentesi).

\begin{soluzione}[: risposte commentate]
\begin{enumerate}[label=\textbf{\arabic*.}]
  \item \textbf{B} (M1).
  \item \textbf{A} (M2). La \texttt{x} interna è un'altra variabile che \q{nasconde} quella
        esterna solo dentro il blocco; fuori, \texttt{x} vale ancora 5.
  \item \textbf{C} (M3). Aritmetica dei puntatori: \texttt{p + 2} punta ad \texttt{a[2]},
        l'asterisco ne legge il valore.
  \item \textbf{B} (M3). La A è un dangling pointer.
  \item \textbf{C} (M4). Regola dei Tre.
  \item \textbf{B} (M4). Gli oggetti locali vengono distrutti in ordine inverso di creazione.
  \item \textbf{B} (M5).
  \item \textbf{D} (M6). Si tiene il termine che cresce di più: $n^2$ cresce più di
        $n \log n$.
  \item \textbf{B} (M6). Il ciclo interno fa $0 + 1 + \ldots + (n-1)$ giri, circa
        $n^2/2$.
  \item \textbf{A} (M7).
  \item \textbf{C} (M7). Il quick sort è $O(n^2)$ nel caso peggiore.
  \item \textbf{C} (M8).
  \item \textbf{B} (M9). Senza \texttt{virtual} il binding è statico: decide il tipo del
        puntatore, \texttt{Base*}.
  \item \textbf{A} (M9). Può però avere costruttori, attributi, ed essere tipo di puntatori.
  \item \textbf{B} (M10). In \texttt{cout << x} l'operando sinistro è \texttt{cout}: un
        metodo dovrebbe appartenere a \texttt{ostream}.
  \item \textbf{A} (M11).
  \item \textbf{B} (M12). Dopo i tre push la pila è 1 2 3 (cima 3): pop dà 3, pop dà 2;
        push(4) mette 4 in cima, pop dà 4.
  \item \textbf{D} (M13).
  \item \textbf{B} (M14).
  \item \textbf{C} (M10). Così si possono concatenare gli assegnamenti
        (\texttt{a = b = c}) senza copie inutili.
\end{enumerate}
\end{soluzione}

\section*{Prova 2 -- Laboratorio: una soluzione completa}

La soluzione segue i punti della traccia con commenti \texttt{// ===== STEP n =====}. Non
è l'unica possibile: conta che ogni punto funzioni e sia mostrato nell'output.
\codice{codice/esame/lab.cpp}
\uscita{codice/esame/output/lab.txt}

\begin{consiglio}[: cosa notare nella soluzione]
\begin{itemize}
  \item \texttt{typeid(...).name()} con MinGW stampa il nome \q{decorato} (\texttt{4Auto}:
        la cifra è la lunghezza del nome). Su altri compilatori può essere diverso: è
        normale.
  \item Coda e BST contengono \texttt{Veicolo*}: i loro distruttori liberano i
        \textbf{nodi}, non i veicoli. I veicoli li ha creati il \texttt{main}, ed è il
        \texttt{main} a liberarli alla fine. Se anche le strutture li cancellassero, ogni
        veicolo verrebbe cancellato due volte.
  \item Nel BST i tre tagliandi da 160 euro hanno la stessa chiave: vanno a destra (regola
        dei duplicati del Modulo 13) e la visita in-order li stampa nell'ordine di
        inserimento.
  \item L'eccezione viene lanciata dentro \texttt{dequeue} e catturata nel \texttt{main}
        \textbf{per riferimento costante}; l'incasso calcolato fino a quel momento resta
        valido.
  \item \texttt{typeid} confronta il tipo \emph{esatto}, \texttt{dynamic\_cast} permette
        anche di usare i metodi propri della derivata (\texttt{getCilindrata}).
  \item Le scelte di progetto discutibili (l'eccezione \texttt{CodaVuota} riusata in
        \texttt{massimo} del BST, le statistiche calcolate nel \texttt{main}) vanno sapute
        giustificare all'orale: \q{per semplicità} è una risposta accettabile, se sai dire
        come faresti meglio (un'eccezione \texttt{AlberoVuoto}, un metodo dedicato).
\end{itemize}
\end{consiglio}

\section*{Prova 3 -- Orale}

Le risposte modello sono nei moduli indicati nella tabella della simulazione. Valuta te
stesso così: hai risposto per almeno un minuto senza fermarti? Hai usato i termini
giusti (glossario)? Hai fatto un esempio o un disegno? Se una delle tre manca, ripeti la
domanda il giorno dopo.
